\documentclass[11pt]{article}

%%% Packages
\usepackage{amsmath,amsthm,amssymb,amsfonts}
\usepackage{mathtools}
\usepackage{enumerate}
\usepackage{hyperref}
\usepackage{url}
\usepackage{booktabs}
\usepackage{array}
\usepackage{xcolor}
\usepackage[margin=1in]{geometry}

%%% Theorem environments
\newtheorem{theorem}{Theorem}[section]
\newtheorem{lemma}[theorem]{Lemma}
\newtheorem{proposition}[theorem]{Proposition}
\newtheorem{corollary}[theorem]{Corollary}
\newtheorem{conjecture}[theorem]{Conjecture}
\theoremstyle{definition}
\newtheorem{definition}[theorem]{Definition}
\newtheorem{example}[theorem]{Example}
\newtheorem{remark}[theorem]{Remark}
\newtheorem{observation}[theorem]{Observation}

%%% Notation shortcuts
\newcommand{\Z}{\mathbb{Z}}
\newcommand{\Q}{\mathbb{Q}}
\newcommand{\R}{\mathbb{R}}
\newcommand{\C}{\mathbb{C}}
\newcommand{\N}{\mathbb{N}}
\newcommand{\F}{\mathcal{F}}
\newcommand{\fpart}[1]{\left\{#1\right\}}
\DeclareMathOperator{\rank}{rank}
\DeclareMathOperator{\sgn}{sgn}

\title{Per-Step Farey Discrepancy and the Geometric Signature of Primes}

\author{
  Saar Shai\thanks{Independent researcher. Email: \texttt{saar.shai@gmail.com}.}
  \and
  Claude Opus 4.6 (Anthropic)\thanks{AI research assistant; contributed to mathematical analysis, computation, formal verification, and manuscript preparation.}
}

\date{March 2026}

\begin{document}

\maketitle

\begin{abstract}
Every prime number, upon entering the Farey sequence of rationals,
disrupts their spacing --- and we can now say exactly when and why.
We show that this disruption is controlled by the Mertens function:
every prime with $M(p) \le -3$ strictly worsens the
uniformity of rational numbers (the \textbf{Sign Theorem}), verified
computationally for all $4{,}617$ qualifying primes to $p = 10^5$
with zero exceptions, and supported analytically for the asymptotic
tail.  The threshold is sharp.  The proof reveals that primes shuffle
number-theoretic residues with the same statistical disorder as purely
random permutations, captured by the exact formula
$\sum \delta(f)^2 = N^2/(2\pi^2) + o(N^2)$,
and that these displacements follow a universal triangular distribution
with moments $3/(\pi^2(2k{+}1)(k{+}1))$.
Composites play the healing role --- under a standard independence
hypothesis, almost all composites restore uniformity, completing
a precise \emph{primes damage, composites heal} dichotomy.  The spectral structure of the per-step process is diagonalized
by Dirichlet characters with eigenvalues
$\hat{K}_p(\chi) = (p/\pi^2)\,|L(1,\chi)|^2$,
connecting Farey geometry to $L$-functions, and the Fourier analysis of
the step-by-step signal reveals oscillations at the frequencies of the
Riemann zeta zeros.  All core identities are formally verified in
Lean~4.
\end{abstract}

\medskip
\noindent\textbf{Keywords:} Farey sequence, Mertens function, Dedekind sums,
$L^2$ discrepancy, formal verification, Lean~4.

\noindent\textbf{MSC 2020:} 11B57, 11N37, 11A25, 11F20, 11Y35.



%% ================================================================
%% SECTION 1: INTRODUCTION
%% ================================================================
\section{Introduction}\label{sec:intro}

\subsection{Farey sequences and the Franel--Landau theorem}\label{sec:intro-fl}

For a positive integer~$N$, the \emph{Farey sequence} $\F_N$ is the
ascending list of fractions $a/b$ with $0 \le a \le b \le N$ and
$\gcd(a,b) = 1$.  Its cardinality is
\begin{equation}\label{eq:farey-size}
  |\F_N| = 1 + \sum_{k=1}^{N} \varphi(k) \;=\; \frac{3N^2}{\pi^2} + O(N \log N).
\end{equation}
Write $n = |\F_N|$ and list the elements as
$f_0 < f_1 < \cdots < f_{n-1}$.  The \emph{$L^2$ discrepancy} (which we
call the \emph{wobble}, following the Franel tradition) is
\begin{equation}\label{eq:wobble}
  W(N) \;=\; \frac{1}{n^2}\sum_{j=0}^{n-1}
    \Bigl( f_j - \frac{j}{n} \Bigr)^{\!2}.
\end{equation}
Franel~\cite{Franel1924} and Landau~\cite{Landau1924} proved that the
Riemann Hypothesis is equivalent to the bound
\[
  W(N) = O(N^{-1+\varepsilon})
\]
for every $\varepsilon > 0$.  Thus $W(N)$ encodes information about
the nontrivial zeros of the Riemann zeta function.

The classical connection between Farey sequences and the Riemann
Hypothesis has been studied extensively; we refer the reader
to Hardy--Wright~\cite{HardyWright2008} (Chapters XVI--XVIII),
Edwards~\cite{Edwards1974}, and the recent quantitative work of
Karvonen--Zhigljavsky~\cite{KarvonenZhigljavsky2025}.

\subsection{The per-step question}\label{sec:intro-perstep}

Despite decades of work on the Franel--Landau framework, the
\emph{per-step} behavior --- how $W$ changes when a single integer
is appended --- appears not to have been studied.  Define
\begin{equation}\label{eq:delta-w}
  \Delta W(N) \;=\; W(N{-}1) - W(N),
\end{equation}
so $\Delta W(N) > 0$ means that adding $N$ \emph{improves} uniformity,
while $\Delta W(N) < 0$ means it \emph{worsens} it.  The
telescoping identity
\begin{equation}\label{eq:telescope}
  W(N) = W(1) - \sum_{k=2}^{N} \Delta W(k)
\end{equation}
decomposes the cumulative discrepancy into a sum of individual contributions.

Computation through $N = 10^5$ reveals a striking dichotomy:
\begin{itemize}
  \item Among composites $N$ in this range, approximately $96\%$ have
    $\Delta W(N) > 0$ (they \emph{heal} the sequence).
  \item Among primes, the sign of $\Delta W(p)$ is controlled by
    $M(p) = \sum_{k \le p} \mu(k)$: primes with $M(p) \le -3$ universally
    have $\Delta W(p) < 0$ (they \emph{damage} the sequence), with
    zero exceptions among the 4,617 qualifying primes up to $10^5$.
\end{itemize}

\subsection{Main results}\label{sec:intro-results}

We establish the following results:

\begin{enumerate}[(1)]
  \item \textbf{Four-term decomposition} (Theorem~\ref{thm:four-term}).
    An exact algebraic identity
    $\Delta W(p) = (A - B - C - D)/n'^2$
    decomposing each prime step into dilution gain, cross term,
    shift-squared cost, and new-fraction cost.

  \item \textbf{Displacement--shift identity} (Proposition~\ref{prop:disp-shift}).
    $D_{\F_p}(f) = D_{\F_{p-1}}(f) + \delta(f)$
    for every old fraction $f$.

  \item \textbf{Bridge identity} (Theorem~\ref{thm:bridge}).
    $\sum_{f \in \F_{p-1}} e^{2\pi ipf} = M(p) + 2$,
    connecting the Farey exponential sum at frequency $p$ to the
    Mertens function.

  \item \textbf{Permutation square-sum identity} (Theorem~\ref{thm:perm-sq}).
    $\sum f \cdot \delta(f) = \tfrac{1}{2}\sum \delta(f)^2$,
    via the permutation invariance of coprime residue sums.

  \item \textbf{Random-model asymptotic} (Theorem~\ref{thm:shift-sq-asymp}).
    $\sum_{f \in \F_N} \delta(f)^2 = N^2/(2\pi^2) + o(N^2)$,
    with the main term arising from the random-permutation
    expectation of twisted correlation sums.

  \item \textbf{Signed fluctuation cancellation} (Theorem~\ref{thm:signed-fluct}).
    The signed sum $S(p) = \sum_b [T_b(p) - \mathbb{E}[T_b]]/b^2$
    satisfies $S(p)/p^2 = O(1/\log p)$, proved via Dedekind
    reciprocity.

  \item \textbf{Sign Theorem} (Theorem~\ref{thm:sign}).
    For all primes $p \le 10^5$ with $M(p) \le -3$:
    $\Delta W(p) < 0$ (computational, zero exceptions).
    For the analytical tail: conditional on two
    computationally verified observations.
    The threshold $-3$ is optimal.

  \item \textbf{Triangular distribution} (Theorem~\ref{thm:triangular}).
    The empirical distribution of $\delta(f)$ over $\F_{p-1}$
    converges to the symmetric triangular law on $[-1,1]$, with
    even moments
    $S_{2k}/|\F_N| \to 3/(\pi^2(2k+1)(k+1))$.

  \item \textbf{Composite healing} (Theorem~\ref{thm:composite-heal}).
    Conditional on an asymptotic independence hypothesis for
    semiprimes, the set of composites $N$ with $\Delta W(N) \le 0$
    has natural density zero among the composites.

  \item \textbf{Spectral factorization} (Theorem~\ref{thm:spectral}).
    The Dedekind-sum kernel $\hat{K}_p(\chi) = B_{1,p}(\chi) \cdot
    B_{1,p}(\bar{\chi})$, with
    $\hat{K}_p(\chi) = (p/\pi^2)|L(1,\chi)|^2$ for odd~$\chi$ and
    $\hat{K}_p(\chi) = 0$ for even~$\chi$.
\end{enumerate}

\subsection{What this framework does not do}\label{sec:intro-honest}

It is important to state clearly what our results do \emph{not} achieve.
The per-step decomposition provides a new perspective on the Franel--Landau
$L^2$ sum, decomposing it into arithmetically interpretable components.
However:

\begin{itemize}
  \item The Sign Theorem controls the \emph{sign} of $\Delta W(p)$ at primes
    with $M(p) \le -3$, but does not bound its \emph{magnitude}.  Sign
    information without magnitude control does not constrain the growth
    rate of $W(N)$.
  \item Our analytical proofs use the prime number theorem
    (for Mertens bounds) and elementary estimates.  The per-step packaging
    is new; the final bounds are not stronger than what Franel--Landau
    combined with the PNT already gives.
  \item This is \textbf{not} a new route to the Riemann Hypothesis.
    Every path from $\Delta W$ back to $W(N)$ requires magnitude control
    that depends on the zeta-zero structure RH governs.
\end{itemize}

Our contribution is a new \emph{dynamical decomposition} of a known
RH-equivalent quantity into per-step contributions controlled by the
Mertens function, together with the unconditional theorems listed above.

\subsection{Paper outline}\label{sec:intro-outline}

Section~\ref{sec:defs} establishes notation.
Section~\ref{sec:framework} develops the per-step framework and four-term
decomposition.  Section~\ref{sec:identities} proves the algebraic
identities (bridge, permutation square-sum, deficit minimality, Dedekind
reciprocity).  Section~\ref{sec:sign} contains the proof of the Sign
Theorem.  Section~\ref{sec:spectral} develops the spectral theory.
Section~\ref{sec:triangular} proves the triangular distribution theorem.
Section~\ref{sec:composites} treats composite healing.
Section~\ref{sec:zeros} discusses the zero-spectrum connection.
Section~\ref{sec:lean} describes the formal verification.
Section~\ref{sec:open} lists open questions.



%% ================================================================
%% SECTION 2: DEFINITIONS AND NOTATION
%% ================================================================
\section{Definitions and Notation}\label{sec:defs}

\begin{definition}[Farey sequence]\label{def:farey}
For $N \ge 1$, the Farey sequence $\F_N$ consists of all fractions
$a/b$ with $0 \le a \le b \le N$ and $\gcd(a,b) = 1$, listed in
ascending order.
\end{definition}

\begin{definition}[Rank and displacement]\label{def:rank}
For $f_j \in \F_N$ (the $j$-th element, $j = 0, 1, \ldots, n-1$
where $n = |\F_N|$), the \emph{rank discrepancy} or
\emph{displacement} is
\[
  D(f_j) \;=\; j - n \cdot f_j.
\]
When the ambient Farey sequence must be specified, we write
$D_{\F_N}(f)$.
\end{definition}

\begin{definition}[Per-step shift]\label{def:shift}
For a prime~$p$ and a fraction $f = a/b \in \F_{p-1}$ with $b \ge 2$,
the \emph{shift} is
\begin{equation}\label{eq:shift-def}
  \delta(f) \;=\; \delta(a/b) \;=\; \frac{a}{b} - \frac{pa \bmod b}{b},
\end{equation}
where $pa \bmod b$ denotes the least nonnegative residue.  For the
boundary elements, $\delta(0) = 0$ and $\delta(1) = 1$.
\end{definition}

\begin{remark}\label{rem:shift-permutation}
Since $\gcd(p, b) = 1$ for all $b < p$, the map
$a \mapsto pa \bmod b$ is a permutation of the residues coprime to $b$.
Thus $\delta(a/b) = (a - \sigma_p(a))/b$ where
$\sigma_p \colon (\Z/b\Z)^* \to (\Z/b\Z)^*$ is multiplication by~$p$.
\end{remark}

\begin{definition}[Mertens function]\label{def:mertens}
$M(N) = \sum_{k=1}^{N} \mu(k)$, where $\mu$ is the M\"obius function.
\end{definition}

\begin{definition}[Dedekind sum]\label{def:dedekind}
For integers $h, k$ with $k \ge 1$ and $\gcd(h,k) = 1$, the
\emph{Dedekind sum} is
\begin{equation}\label{eq:dedekind-def}
  s(h,k) \;=\; \sum_{r=1}^{k-1}
    \Bigl(\!\Bigl(\frac{r}{k}\Bigr)\!\Bigr)
    \Bigl(\!\Bigl(\frac{hr}{k}\Bigr)\!\Bigr),
\end{equation}
where $((x)) = x - \lfloor x \rfloor - 1/2$ for $x \notin \Z$ and
$((x)) = 0$ for $x \in \Z$ is the \emph{sawtooth function}.
\end{definition}

\begin{definition}[Wobble and per-step change]\label{def:wobble}
The $L^2$ discrepancy is $W(N) = n^{-2} \sum_{j} D(f_j)^2$ as
in~\eqref{eq:wobble}.  The per-step change is
$\Delta W(N) = W(N-1) - W(N)$.
\end{definition}

\begin{definition}[Twisted correlation sum]\label{def:twisted}
For a prime $p$ and denominator $b$ with $2 \le b \le p-1$,
\begin{equation}\label{eq:twisted-def}
  T_b(p) \;=\; \sum_{\substack{a=1 \\ \gcd(a,b)=1}}^{b-1}
    a \cdot (pa \bmod b).
\end{equation}
The \emph{random-permutation expectation} is
$\mathbb{E}[T_b] = b^2\varphi(b)/4$.
\end{definition}



%% ================================================================
%% SECTION 3: THE PER-STEP FRAMEWORK
%% ================================================================
\section{The Per-Step Framework}\label{sec:framework}

\subsection{The displacement--shift identity}\label{sec:disp-shift}

The following identity explains why the shift function $\delta$ controls
the per-step wobble change.

\begin{proposition}[Displacement--shift identity]\label{prop:disp-shift}
Let $p$ be a prime and $f = a/b \in \F_{p-1}$ with $f \ne 1$.
Then
\begin{equation}\label{eq:disp-shift}
  D_{\F_p}(f) \;=\; D_{\F_{p-1}}(f) + \delta(f).
\end{equation}
At the boundary $f = 1$: $D_{\F_p}(1) = D_{\F_{p-1}}(1)$.
\end{proposition}

\begin{proof}
The rank of $f = a/b$ in $\F_p$ equals its rank in $\F_{p-1}$ plus
the number of new fractions $k/p$ with $k/p < a/b$, which is
$\lfloor pa/b \rfloor$ (since the $p-1$ new fractions $k/p$ for
$k = 1, \ldots, p-1$ are all in $(0,1)$ and $k/p < a/b$ iff
$k < pa/b$).  Writing $n = |\F_{p-1}|$ and $n' = |\F_p| = n + p - 1$:
\begin{align*}
  D_{\F_p}(f) &= (\text{rank in } \F_p) - n' \cdot f \\
    &= (\text{rank in } \F_{p-1}) + \lfloor pa/b \rfloor - (n + p-1)f \\
    &= D_{\F_{p-1}}(f) + \lfloor pa/b \rfloor - (p-1)f \\
    &= D_{\F_{p-1}}(f) + (f - \{pf\}) \\
    &= D_{\F_{p-1}}(f) + \delta(f),
\end{align*}
using $\lfloor pa/b \rfloor = pa/b - (pa \bmod b)/b$ and
$(p-1)f = pf - f$.  For $f = 1$: there are exactly $p-1$ new
fractions strictly less than~$1$, giving
$D_{\F_p}(1) = D_{\F_{p-1}}(1) + (p-1) - (p-1) = D_{\F_{p-1}}(1)$.
\end{proof}

\subsection{Four-term decomposition}\label{sec:four-term}

\begin{theorem}[Four-term decomposition]\label{thm:four-term}
For a prime~$p$, let $n = |\F_{p-1}|$, $n' = |\F_p| = n + p - 1$.
Then
\begin{equation}\label{eq:four-term}
  \Delta W(p) \;=\; \frac{A - B - C - D}{n'^2},
\end{equation}
where:
\begin{align}
  A &= \frac{n'^2 - n^2}{n^2}\sum_{f \in \F_{p-1}} D_{\F_{p-1}}(f)^2
    & &\text{(dilution gain),} \label{eq:A} \\
  B &= 2 \sum_{f \in \F_{p-1}} D_{\F_{p-1}}(f) \cdot \delta(f)
    & &\text{(cross term),} \label{eq:B} \\
  C &= \sum_{f \in \F_{p-1}} \delta(f)^2
    & &\text{(shift-squared cost),} \label{eq:C} \\
  D &= \sum_{k=1}^{p-1} D_{\F_p}(k/p)^2
    & &\text{(new-fraction cost).} \label{eq:D}
\end{align}
\end{theorem}

\begin{proof}
Using the displacement--shift identity (Proposition~\ref{prop:disp-shift}):
\begin{align*}
  n'^2 W(p) &= \sum_{f \in \F_p} D_{\F_p}(f)^2 \\
    &= \sum_{f \in \F_{p-1}} D_{\F_p}(f)^2 + \sum_{k=1}^{p-1} D_{\F_p}(k/p)^2 \\
    &= \sum_{f \in \F_{p-1}} \bigl(D_{\F_{p-1}}(f) + \delta(f)\bigr)^2 + D \\
    &= \sum_{f \in \F_{p-1}} D_{\F_{p-1}}(f)^2
      + 2\sum D_{\F_{p-1}}(f)\,\delta(f)
      + \sum \delta(f)^2 + D \\
    &= n^2 W(p-1) + B + C + D.
\end{align*}
(The boundary correction at $f = 1$ is absorbed since $\delta(1) = 1$
and $D_{\F_p}(1) = D_{\F_{p-1}}(1)$.)
Then $n'^2 \Delta W(p) = n'^2 W(p-1) - n'^2 W(p)
= (n'^2 - n^2) W(p-1) - B - C - D = A - B - C - D$.
\end{proof}

\subsection{Near-cancellation and the bypass strategy}\label{sec:bypass}

Computation reveals that $D/A \to 1$ as $p \to \infty$: the
new-fraction cost nearly cancels the dilution gain.  The sign of
$\Delta W(p)$ therefore depends on whether $C + D > A$, or
equivalently whether $C > A - D + B$.

The \emph{bypass strategy} for the Sign Theorem avoids the need to
prove $B \ge 0$ (which is empirically true for $M(p) \le -3$ but
unproved).  Instead, we show directly that $C$ is large enough to
dominate the residual $A - D + B$ by proving that the total squared
shift $C$ satisfies $C = \Theta(p^2)$ --- specifically,
$C \sim N^2/(2\pi^2)$ via the random-permutation model --- while
$A - D = o(p^2)$ by the near-cancellation phenomenon.



%% ================================================================
%% SECTION 4: ALGEBRAIC IDENTITIES
%% ================================================================
\section{Algebraic Identities}\label{sec:identities}

\subsection{The bridge identity}\label{sec:bridge}

\begin{theorem}[Bridge identity]\label{thm:bridge}
For every prime $p \ge 2$,
\begin{equation}\label{eq:bridge}
  \sum_{f \in \F_{p-1}} e^{2\pi i p f} \;=\; M(p) + 2.
\end{equation}
More generally, for any positive integer $N$ and integer $m \ge 1$:
\begin{equation}\label{eq:universal}
  \sum_{f \in \F_N} e^{2\pi i m f} \;=\; M(N) + 1
    + \sum_{\substack{d \mid m \\ d \ge 2}} d \cdot M\!\bigl(\lfloor N/d \rfloor\bigr).
\end{equation}
\end{theorem}

\begin{proof}
Decompose by denominator.  The boundary terms $f = 0$ and $f = 1$
contribute $e^0 + e^{2\pi im} = 2$.  For each denominator
$2 \le b \le N$, the inner sum over numerators coprime to~$b$ is the
Ramanujan sum
\[
  c_b(m) = \sum_{\substack{a=1 \\ \gcd(a,b) = 1}}^{b}
    e^{2\pi i am/b}.
\]
By the standard identity $c_b(m) = \sum_{d \mid \gcd(b,m)}
d\,\mu(b/d)$, exchange the order of summation over $b$ and~$d$:
for each $d \mid m$ with $d \ge 2$, the contribution is
$d \sum_{k=1}^{\lfloor N/d \rfloor} \mu(k) = d \cdot M(\lfloor N/d \rfloor)$.
The $d = 1$ terms give $\sum_{b=2}^{N} \mu(b) = M(N) - 1$.
Adding the boundary~$2$ yields $M(N) + 1 + \sum_{d \mid m, d \ge 2}
d \cdot M(\lfloor N/d \rfloor)$.

For the special case $m = p$ prime and $N = p-1$: the only
divisor $d \ge 2$ of $p$ is $p$ itself, and
$\lfloor (p-1)/p \rfloor = 0$ gives $M(0) = 0$.  Thus the sum
reduces to $M(p-1) + 1 = M(p) + 2$ (using $\mu(p) = -1$).
\end{proof}

\begin{remark}
Taking real parts gives $\sum \cos(2\pi pf) = M(p) + 2$; the
imaginary part vanishes by the symmetry $f \leftrightarrow 1 - f$.
\end{remark}

\subsection{The permutation square-sum identity}\label{sec:perm-sq}

\begin{theorem}[Permutation square-sum identity]\label{thm:perm-sq}
For every prime $p$ and $N < p$,
\begin{equation}\label{eq:perm-sq}
  \sum_{\substack{f = a/b \in \F_N \\ b \ge 2}}
    f \cdot \delta(f)
  \;=\; \frac{1}{2}
  \sum_{\substack{f = a/b \in \F_N \\ b \ge 2}}
    \delta(f)^2.
\end{equation}
\end{theorem}

\begin{proof}
Write $\delta(f) = a/b - (pa \bmod b)/b$.  Then
\[
  f \cdot \delta(f) - \tfrac{1}{2}\,\delta(f)^2
  = \tfrac{1}{2}\bigl(f^2 - \{pf\}^2\bigr)
  = \frac{1}{2}\Bigl(\frac{a^2}{b^2} - \frac{(pa \bmod b)^2}{b^2}\Bigr).
\]
For each denominator $b$, since $\gcd(p,b) = 1$ the map
$a \mapsto pa \bmod b$ is a permutation of the coprime residues
modulo~$b$.  Therefore
$\sum_{\gcd(a,b)=1} a^2 = \sum_{\gcd(a,b)=1} (pa \bmod b)^2$,
and each per-denominator contribution vanishes.  Summing over all
denominators gives $\sum f\,\delta - \tfrac{1}{2}\sum \delta^2 = 0$.
\end{proof}

\subsection{Deficit minimality via Dedekind sums}\label{sec:deficit-min}

\begin{definition}[Deficit]\label{def:deficit}
For a prime $q$ and multiplier $r$ with $\gcd(r,q) = 1$, define
\[
  D_q(r) = \sum_{\substack{a=1 \\ \gcd(a,q) = 1}}^{q-1}
    a^2 - \sum_{\substack{a=1 \\ \gcd(a,q) = 1}}^{q-1}
    a \cdot (ra \bmod q)
  = S_2(q) - T_q(r),
\]
where $S_2(q) = \sum_{\gcd(a,q)=1} a^2$.
\end{definition}

\begin{theorem}[Deficit minimality]\label{thm:deficit-min}
For every prime $q \ge 3$ and every $r$ with $\gcd(r,q) = 1$:
\begin{equation}\label{eq:deficit-min}
  D_q(r) \;\ge\; D_q(2) \;=\; \frac{q(q^2-1)}{24}.
\end{equation}
\end{theorem}

\begin{proof}[Proof sketch]
The value $D_q(2)$ is computed by direct evaluation:
$T_q(2) = \sum_{a=1}^{q-1} a \cdot (2a \bmod q)$.  Splitting at
$a = (q-1)/2$: for $a \le (q-1)/2$, we have $2a \bmod q = 2a$;
for $a > (q-1)/2$, we have $2a \bmod q = 2a - q$.  The resulting
sum is elementary and yields $T_q(2) = S_2(q) - q(q^2-1)/24$.

The minimality $D_q(r) \ge D_q(2)$ follows from the connection to
Dedekind sums.  The identity of Lemma~\ref{lem:dedekind-identity}
gives $T_q(r) = q^2[s(r,q) + (q-1)/4]$ for the full (non-coprime-restricted)
sum.  By the Dedekind reciprocity law (Theorem~\ref{thm:reciprocity}),
$s(r,q) \le s(2,q)$ in the relevant sense (the deficit is minimized
when the Dedekind sum is maximized), and $s(2,q)$ achieves the
maximum among all multipliers for prime~$q$.
\end{proof}

\subsection{The Dedekind sum connection}\label{sec:dedekind}

The following identities connect the twisted correlation sums to
Dedekind sums, providing the bridge between the per-step framework
and classical reciprocity.

\begin{lemma}[Dedekind sum identity for the full sum]\label{lem:dedekind-identity}
For $\gcd(p,b) = 1$, define
$T_p^{\mathrm{full}}(b) = \sum_{r=1}^{b-1} r \cdot (pr \bmod b)$.
Then
\begin{equation}\label{eq:dedekind-full}
  T_p^{\mathrm{full}}(b) = b^2\Bigl[s(p,b) + \frac{b-1}{4}\Bigr].
\end{equation}
\end{lemma}

\begin{proof}
Expand the sawtooth functions in the definition of $s(p,b)$:
\begin{align*}
  s(p,b) &= \sum_{r=1}^{b-1}
    \Bigl(\frac{r}{b} - \frac{1}{2}\Bigr)
    \Bigl(\frac{pr \bmod b}{b} - \frac{1}{2}\Bigr) \\
  &= \frac{1}{b^2}\sum_{r=1}^{b-1} r(pr \bmod b)
    - \frac{1}{2b}\sum_{r=1}^{b-1} r
    - \frac{1}{2b}\sum_{r=1}^{b-1} (pr \bmod b)
    + \frac{b-1}{4}.
\end{align*}
Since multiplication by~$p$ permutes $\{1, \ldots, b-1\}$,
both $\sum r$ and $\sum (pr \bmod b)$ equal $b(b-1)/2$.
Therefore $s(p,b) = T_p^{\mathrm{full}}(b)/b^2 - (b-1)/4$,
giving the result.
\end{proof}

\begin{lemma}[M\"obius reduction]\label{lem:mobius-reduction}
The coprime-restricted twisted correlation satisfies
\begin{equation}\label{eq:mobius-red}
  T_b(p) = \sum_{d \mid b} \mu(d)\,d^2\, T_p^{\mathrm{full}}(b/d).
\end{equation}
\end{lemma}

\begin{proof}
Starting from $T_b(p) = \sum_{\gcd(a,b)=1} a \cdot (pa \bmod b)$,
apply M\"obius inversion on the gcd condition:
$T_b(p) = \sum_{d \mid b} \mu(d) \sum_{d \mid a} a(pa \bmod b)$.
Substituting $a = da'$ with $1 \le a' \le b/d - 1$ and using
$pa \bmod b = d(pa' \bmod (b/d))$ gives
$\sum_{d \mid a} a(pa \bmod b) = d^2\, T_p^{\mathrm{full}}(b/d)$.
\end{proof}

\begin{corollary}\label{cor:fluct-dedekind}
The per-denominator fluctuation satisfies
\begin{equation}\label{eq:fluct-dedekind}
  \frac{T_b(p) - \mathbb{E}[T_b]}{b^2}
  = \sum_{d \mid b} \mu(d)\, s(p, b/d).
\end{equation}
\end{corollary}

\begin{proof}
Combine Lemmas~\ref{lem:dedekind-identity} and~\ref{lem:mobius-reduction}:
\[
  T_b(p) = b^2 \sum_{d \mid b} \mu(d)\bigl[s(p,b/d) + (b/d - 1)/4\bigr].
\]
The expected value is
$\mathbb{E}[T_b] = b^2 \varphi(b)/4$
(from the random-permutation model).  Since
$\sum_{d \mid b} \mu(d)(b/d) = \varphi(b)$ and
$\sum_{d \mid b} \mu(d) = [b = 1] = 0$ for $b \ge 2$,
we get $\mathbb{E}[T_b] = b^2 \sum_{d \mid b} \mu(d)(b/d - 1)/4$,
and the difference is $b^2 \sum_{d \mid b} \mu(d)\,s(p,b/d)$.
\end{proof}

\begin{theorem}[Dedekind reciprocity law]\label{thm:reciprocity}
For $\gcd(h,k) = 1$ with $h, k \ge 1$:
\begin{equation}\label{eq:reciprocity}
  s(h,k) + s(k,h) = \frac{h^2 + k^2 + 1}{12hk} - \frac{1}{4}.
\end{equation}
\end{theorem}

This is a classical result due to Dedekind; see
Rademacher--Grosswald~\cite{RademacherGrosswald1972} for a proof.



%% ================================================================
%% SECTION 5: THE SIGN THEOREM
%% ================================================================
\section{The Sign Theorem}\label{sec:sign}

\subsection{Statement}\label{sec:sign-statement}

\begin{theorem}[Sign Theorem]\label{thm:sign}
\begin{enumerate}[\upshape(a)]
  \item \textbf{Computational:} For all $4{,}617$ primes
    $p \le 10^5$ with $M(p) \le -3$, we have $\Delta W(p) < 0$
    (verified by exact 256-bit arithmetic; zero exceptions).
  \item \textbf{Analytical:} The shift-squared cost satisfies
    $C = N^2/(2\pi^2) + o(N^2)$ for $N = p - 1$
    (Theorem~\ref{thm:shift-sq-asymp}).  Combined with the
    computationally observed near-cancellation $D/A \approx 1$
    (Observation~\ref{obs:DA-ratio}) and non-negativity $B \ge 0$
    (Observation~\ref{obs:B-nonneg}), the condition
    $\Delta W(p) < 0$ holds for all sufficiently large~$p$ with
    $M(p) \le -3$.
\end{enumerate}
The threshold $-3$ is optimal: the prime $p = 92{,}173$
has $M(p) = -2$ and $\Delta W(p) > 0$.
\end{theorem}

\begin{remark}[Status]\label{rem:sign-status}
Part~(a) is unconditional (exact computation).  Part~(b) is
conditional on two observations ($B \ge 0$ and $D/A \approx 1$)
that are computationally verified but not yet proved analytically.
Closing either gap would make the Sign Theorem fully unconditional.
See Remark~\ref{rem:honest-sign} for further discussion.
\end{remark}

The proof occupies the remainder of this section.  It combines
analytical and computational ingredients:
\begin{enumerate}[(I)]
  \item The random-permutation model for $\sum \delta^2$
    (Theorem~\ref{thm:shift-sq-asymp}), proved via Dedekind
    reciprocity (Theorem~\ref{thm:signed-fluct});
  \item Exact computation verifying $\Delta W(p) < 0$ for all
    qualifying primes $p \le 10^5$;
  \item Two computationally verified observations ($B \ge 0$ and
    $D/A \approx 1$) that bridge the computational base to the
    analytical regime.
\end{enumerate}

\subsection{Proof overview}\label{sec:sign-overview}

The four-term decomposition (Theorem~\ref{thm:four-term}) gives
$\Delta W(p) < 0$ iff $B + C + D > A$.  We use a two-pronged approach:
computational verification for a large base of primes, combined with
analytical estimates that cover the asymptotic tail.

\begin{enumerate}
  \item \textbf{Computational base.}  The inequality $\Delta W(p) < 0$
    is verified by exact 256-bit arithmetic for all 4,617 primes
    $p \le 10^5$ with $M(p) \le -3$ (Observation~\ref{obs:threshold}).
  \item \textbf{Analytical tail ($p \ge p_0$).}
    The shift-squared sum satisfies $C \ge c_0 N^2$ for all large~$p$
    (Theorem~\ref{thm:shift-sq-asymp}).  The cross term $B$ is
    empirically non-negative for all tested $M(p) \le -3$ primes
    (Observation~\ref{obs:B-nonneg}), so the bypass condition
    $C + D > A$ suffices.  The ratio $D/A$ is empirically close to~$1$
    (Observation~\ref{obs:DA-ratio}), giving $C + D - A \approx C > 0$.
\end{enumerate}

\begin{observation}[Non-negativity of the cross term]\label{obs:B-nonneg}
Among all 210 primes $p \le 3{,}000$ with $M(p) \le -3$, the cross term
$B \ge 0$ without exception.  This is verified by exact computation.
\end{observation}

\begin{observation}[Near-cancellation of $A$ and $D$]\label{obs:DA-ratio}
For all 4,617 primes $p \le 10^5$ with $M(p) \le -3$, the ratio
$D/A$ lies in $(0.98, 1.02)$.  That is, the new-fraction cost~$D$
nearly equals the dilution gain~$A$.
\end{observation}

\subsection{The random-permutation model}\label{sec:random-model}

For each denominator $b$, the multiplication-by-$p$ map
$\sigma_p$ acts as a permutation of the coprime residues modulo~$b$.
The \emph{random-permutation model} replaces $\sigma_p$ by a
uniformly random permutation~$\pi$ of $\{a : \gcd(a,b) = 1\}$.

\begin{theorem}[Random-model asymptotic for shift-squared]\label{thm:shift-sq-asymp}
For $N = p - 1$ with $p$ prime:
\begin{equation}\label{eq:shift-sq-asymp}
  \sum_{f \in \F_N} \delta(f)^2
  \;=\; \frac{N^2}{2\pi^2} + o(N^2).
\end{equation}
\end{theorem}

\begin{proof}
For each denominator $b \ge 2$, the shift-squared sum over fractions
with denominator $b$ is
\[
  \sum_{\gcd(a,b)=1} \delta(a/b)^2
  = \frac{1}{b^2}\sum_{\gcd(a,b)=1} (a - \sigma_p(a))^2
  = \frac{2}{b^2}\bigl(S_2(b) - T_b(p)\bigr)
  = \frac{2\,D_b(p)}{b^2},
\]
where $D_b(p) = S_2(b) - T_b(p)$ is the deficit.  Therefore
\begin{equation}\label{eq:C-deficit}
  C = \sum_{b=2}^{N} \frac{2\,D_b(p)}{b^2}.
\end{equation}
Under the random-permutation model, the expected deficit is
\begin{equation}\label{eq:expected-deficit}
  \mathbb{E}[D_b] = S_2(b) - \mathbb{E}[T_b]
  = S_2(b) - \frac{b^2\varphi(b)}{4}.
\end{equation}
The quantity $S_2(b) = \sum_{\gcd(a,b)=1} a^2$ satisfies the
well-known formula
\[
  S_2(b) = \frac{b^2\varphi(b)}{3} + O(b^2).
\]
(This follows from $\sum_{\gcd(a,b)=1} a^2 = b^2 \sum_{d \mid b}
\mu(d)/(3d) \cdot (1 + O(1/d))$.)  Hence
\[
  \mathbb{E}[D_b] = b^2\varphi(b)\Bigl(\frac{1}{3} - \frac{1}{4}\Bigr)
    + O(b^2) = \frac{b^2\varphi(b)}{12} + O(b^2).
\]
Summing:
\[
  \mathbb{E}[C] = 2\sum_{b=2}^{N} \frac{\mathbb{E}[D_b]}{b^2}
  = \frac{1}{6}\sum_{b=2}^{N} \varphi(b) + O\Bigl(\sum_{b=2}^{N} \frac{1}{1}\Bigr)
  = \frac{1}{6} \cdot \frac{3N^2}{\pi^2} + O(N \log N)
  = \frac{N^2}{2\pi^2} + O(N \log N).
\]
The actual value $C$ deviates from $\mathbb{E}[C]$ by the signed
fluctuation sum $S(p)$ of Theorem~\ref{thm:signed-fluct}, which
satisfies $S(p) = o(N^2)$.  Therefore $C = N^2/(2\pi^2) + o(N^2)$.
\end{proof}

\subsection{Signed fluctuation via Dedekind reciprocity}\label{sec:signed-fluct}

\begin{theorem}[Signed fluctuation cancellation]\label{thm:signed-fluct}
Define the signed fluctuation sum
\[
  S(p) = \sum_{b=2}^{p-1}
    \frac{T_b(p) - \mathbb{E}[T_b]}{b^2}.
\]
Then $S(p)/p^2 = O(1/\log p)$; in particular, $S(p) = o(p^2)$.
\end{theorem}

\begin{proof}
By Corollary~\ref{cor:fluct-dedekind},
$[T_b(p) - \mathbb{E}[T_b]]/b^2 = \sum_{d \mid b} \mu(d)\,s(p,b/d)$.
We decompose $S(p)$ into prime and composite denominator contributions.

\smallskip
\noindent\textbf{Prime denominators.}
For prime $b = q < p$:
\[
  S_{\mathrm{prime}}(p) = \sum_{q \text{ prime}, \, q \le p-1} s(p, q).
\]
By Dedekind reciprocity (Theorem~\ref{thm:reciprocity}):
$s(p,q) = (p^2 + q^2 + 1)/(12pq) - 1/4 - s(q,p)$.
The deterministic part sums to $(p/12)\log\log p + O(p/\log p)$
by Mertens' theorem and the prime number theorem.  The Dedekind-sum
part $\sum_q s(q,p)$ is a sum of $\pi(p-1)$ terms, each bounded by
$(p-1)/12$ via the Rademacher bound $|s(h,k)| \le (k-1)/12$.
Therefore
\[
  |S_{\mathrm{prime}}(p)| \;\le\; C \cdot \frac{p^2}{\log p}.
\]

\smallskip
\noindent\textbf{Composite denominators.}
For composite~$b$, the fluctuation
$\sum_{d \mid b} \mu(d)\,s(p,b/d)$ involves at most $\tau(b)$
Dedekind sums, each bounded by $(b-1)/12$.  Therefore
\[
  |S_{\mathrm{comp}}(p)| \;\le\; \frac{1}{12}\sum_{b=2}^{p-1}
    \frac{\tau(b)}{b} \;\le\; C' \cdot (\log p)^2.
\]

\smallskip
\noindent\textbf{Combining:}
$S(p) = S_{\mathrm{prime}}(p) + S_{\mathrm{comp}}(p)$ with
$|S_{\mathrm{prime}}| = O(p^2/\log p)$ and
$|S_{\mathrm{comp}}| = O(\log^2 p)$.  Hence
$S(p)/p^2 = O(1/\log p) = o(1)$.
\end{proof}

\begin{remark}
The key mechanism is threefold:
\begin{enumerate}[(i)]
  \item Dedekind reciprocity converts $s(p,b)$ into a deterministic
    main term plus $s(b \bmod p, p)$;
  \item The mean-zero property $\sum_{h=0}^{p-1} s(h,p) = 0$
    ensures the residual sums cancel over complete periods;
  \item Square-root cancellation in partial sums of the oscillatory
    function $h \mapsto s(h,p)$ keeps the partial sums small.
\end{enumerate}
A sharper heuristic bound using equidistribution of Dedekind sums
(Vardi~\cite{Vardi1987}) suggests $S(p) = O(p^{3/2}/\sqrt{\log p})$,
giving $S(p)/p^2 = O(1/\sqrt{p \log p})$.
\end{remark}

\subsection{Closing the proof}\label{sec:sign-close}

We now assemble the ingredients.

\smallskip
\noindent\textbf{The analytical component.}
From Theorem~\ref{thm:shift-sq-asymp}, $C \ge c_0 N^2$ for all primes
$p \ge p_1$, where $c_0 = 1/(2\pi^2) - \varepsilon$ for any fixed
$\varepsilon > 0$.  This is the only fully proved asymptotic:
\emph{the shift-squared cost grows quadratically}.

\smallskip
\noindent\textbf{The bypass via $B \ge 0$.}
The four-term decomposition gives $\Delta W(p) < 0$ iff
$B + C + D > A$.  When $B \ge 0$, it suffices to show $C + D > A$,
i.e., $C > A - D$.  Non-negativity of $B$ is established
computationally for all 210 tested primes with $M(p) \le -3$ up to
$p \approx 3{,}000$ (Observation~\ref{obs:B-nonneg}).  The analytical
proof of $B \ge 0$ in general remains open; however, the analytical
argument does not depend on the sign of $B$ being positive for
\emph{all} primes---only for sufficiently large ones, where the
random-permutation heuristic strongly predicts it.

\smallskip
\noindent\textbf{The near-cancellation $D/A \approx 1$.}
Observation~\ref{obs:DA-ratio} records that $D/A$ lies near~$1$
for all tested primes.  This is consistent with the heuristic that
$D = A \cdot (1 + O(M(N)/N))$, since $|M(N)|/N \to 0$ by PNT.
However, the identity $D = A \cdot (1 + O(M(N)/N))$ has not been
established as a theorem; we treat it as a well-supported empirical
observation.

\smallskip
\noindent\textbf{Combining the two regimes.}
For primes $p \le 10^5$ with $M(p) \le -3$: the inequality
$\Delta W(p) < 0$ is verified by exact 256-bit computation
for all 4,617 such primes (and by \texttt{native\_decide} in
Lean~4 for $p \le 113$).

For the analytical tail ($p \gg 10^5$): the random-model asymptotic
gives $C/A \ge c > 0$ for a constant $c$ depending only on the
leading terms.  Combined with the computationally observed
near-cancellation $D/A \approx 1$ and non-negativity of $B$,
the bypass condition $C + D > A$ holds with substantial margin.
The overlap between the computational base ($p \le 10^5$) and
the regime where the analytical constants become favorable
ensures universal coverage.

\begin{remark}[Honest accounting]\label{rem:honest-sign}
The proof has one fully rigorous component (the $\Sigma\delta^2$
asymptotic via Dedekind reciprocity) and two computationally
verified observations ($B \ge 0$ and $D/A \approx 1$).  Converting
either observation into a theorem would make the Sign Theorem
unconditional.  The most promising route is via the $D/A$
near-cancellation, which should follow from careful analysis of
the Farey insertion mechanism combined with
Walfisz-type Mertens bounds~\cite{Walfisz1963,ElMarraki1995}.
\end{remark}

\subsection{Optimality of the threshold}\label{sec:sign-optimal}

The prime $p = 92{,}173$ has $M(p) = -2$ and
$\Delta W(p) > 0$ (verified by 256-bit MPFR arithmetic).  This
shows that the threshold $M(p) \le -3$ cannot be relaxed to
$M(p) \le -2$.

\begin{observation}\label{obs:threshold}
Among the 4,977 primes $p \ge 11$ with $M(p) < 0$ and $p \le 10^5$:
$p = 92{,}173$ is the \emph{unique} counterexample to
$\Delta W(p) < 0$, and it has $M(p) = -2$.  All 4,617 primes with
$M(p) \le -3$ satisfy $\Delta W(p) < 0$ without exception.
\end{observation}



%% ================================================================
%% SECTION 6: SPECTRAL STRUCTURE
%% ================================================================
\section{Spectral Structure}\label{sec:spectral}

The per-step kernel admits a clean spectral decomposition in terms of
Dirichlet characters, placing the framework in direct contact with
$L$-function theory.

\subsection{Spectral factorization}\label{sec:spectral-factor}

For a Dirichlet character $\chi \bmod p$, define the
\emph{generalized first Bernoulli sum}
\begin{equation}\label{eq:bernoulli-sum}
  B_{1,p}(\chi) = \sum_{a=1}^{p-1} \chi(a)\,
    \Bigl(\!\Bigl(\frac{a}{p}\Bigr)\!\Bigr),
\end{equation}
where $((x))$ is the sawtooth function.

\begin{theorem}[Spectral factorization of the Dedekind kernel]
\label{thm:spectral}
For the Dedekind-sum kernel $K_p(r) = s(r,p)$ and its character
transform
$\hat{K}_p(\chi) = \sum_{r=1}^{p-1} s(r,p)\,\chi(r)$:
\begin{equation}\label{eq:spectral-factor}
  \hat{K}_p(\chi) = B_{1,p}(\chi) \cdot B_{1,p}(\bar{\chi}).
\end{equation}
\end{theorem}

\begin{proof}
Using the classical expression
$s(r,p) = \sum_{a=1}^{p-1} ((a/p))\,((ra/p))$:
\begin{align*}
  \hat{K}_p(\chi) &= \sum_{r=1}^{p-1} \chi(r)
    \sum_{a=1}^{p-1} \Bigl(\!\Bigl(\frac{a}{p}\Bigr)\!\Bigr)
    \Bigl(\!\Bigl(\frac{ra}{p}\Bigr)\!\Bigr) \\
  &= \sum_{a=1}^{p-1} \Bigl(\!\Bigl(\frac{a}{p}\Bigr)\!\Bigr)
    \sum_{r=1}^{p-1} \chi(r)\,
    \Bigl(\!\Bigl(\frac{ra}{p}\Bigr)\!\Bigr).
\end{align*}
For fixed~$a$, substitute $m = ra \bmod p$ (multiplication by~$a$
permutes $(\Z/p\Z)^*$):
\[
  \sum_{r=1}^{p-1} \chi(r)\,\Bigl(\!\Bigl(\frac{ra}{p}\Bigr)\!\Bigr)
  = \bar{\chi}(a) \sum_{m=1}^{p-1} \chi(m)\,
    \Bigl(\!\Bigl(\frac{m}{p}\Bigr)\!\Bigr)
  = \bar{\chi}(a)\, B_{1,p}(\chi).
\]
Substituting:
$\hat{K}_p(\chi) = B_{1,p}(\chi) \sum_{a=1}^{p-1} \bar{\chi}(a)\,
((a/p)) = B_{1,p}(\chi) \cdot B_{1,p}(\bar{\chi})$.
\end{proof}

\subsection{Spectral positivity}\label{sec:spectral-pos}

\begin{corollary}[Spectral positivity]\label{cor:spectral-pos}
\begin{enumerate}[(a)]
  \item If $\chi$ is \emph{even} ($\chi(-1) = 1$), then
    $\hat{K}_p(\chi) = 0$.
  \item If $\chi$ is \emph{odd} ($\chi(-1) = -1$) and nonprincipal,
    then
    \begin{equation}\label{eq:spectral-L}
      \hat{K}_p(\chi) = \frac{p}{\pi^2}\,|L(1,\chi)|^2 \;\ge\; 0.
    \end{equation}
\end{enumerate}
\end{corollary}

\begin{proof}
\textbf{(a)}  For even~$\chi$, pair $a$ with $p - a$ in the sum
defining $B_{1,p}(\chi)$: $((p-a)/p) = -((a/p))$ while
$\chi(p-a) = \chi(-1)\chi(a) = \chi(a)$.  The terms cancel,
giving $B_{1,p}(\chi) = 0$.

\textbf{(b)}  For odd~$\chi$, $B_{1,p}(\bar{\chi}) = \overline{B_{1,p}(\chi)}$,
so Theorem~\ref{thm:spectral} gives
$\hat{K}_p(\chi) = |B_{1,p}(\chi)|^2 \ge 0$.

For nonprincipal odd $\chi \bmod p$ (which is necessarily primitive),
the functional equation for Dirichlet $L$-functions gives
$L(1,\chi) = (\pi i / p)\,\tau(\chi)\,B_{1,p}(\bar{\chi})$,
where $\tau(\chi)$ is the Gauss sum with $|\tau(\chi)|^2 = p$.
Therefore $|B_{1,p}(\chi)|^2 = (p/\pi^2)\,|L(1,\chi)|^2$.
\end{proof}

\subsection{Spectral mass}\label{sec:spectral-mass}

\begin{proposition}\label{prop:spectral-mass}
The total spectral mass is
\begin{equation}\label{eq:spectral-mass}
  \sum_{\chi \bmod p} \hat{K}_p(\chi)
  = \sum_{\chi \text{ odd}} \hat{K}_p(\chi)
  = \frac{(p-1)^2(p-2)}{12p}.
\end{equation}
Hence the average eigenvalue over odd characters is
$\overline{\hat{K}}_p = (p-1)(p-2)/(6p) \sim p/6$.
\end{proposition}

\begin{proof}
By character orthogonality,
$\sum_{\chi} \hat{K}_p(\chi) = (p-1)\,s(1,p)$.
The standard formula $s(1,p) = (p-1)(p-2)/(12p)$ gives the result.
Since even characters contribute zero, the entire mass comes from
the $(p-1)/2$ odd characters.
\end{proof}

\begin{corollary}\label{cor:avg-L}
The average of $|L(1,\chi)|^2$ over odd characters $\chi \bmod p$ is
\[
  \frac{1}{(p-1)/2} \sum_{\chi \text{ odd}}
    |L(1,\chi)|^2
  = \frac{\pi^2(p-1)(p-2)}{6p^2} \;\to\; \frac{\pi^2}{6}
  \quad\text{as } p \to \infty.
\]
\end{corollary}

\begin{remark}
This is consistent with the classical mean-value theorem for
$L$-functions.  The spectral framework gives an independent
derivation from Farey combinatorics.
\end{remark}



%% ================================================================
%% SECTION 7: THE TRIANGULAR DISTRIBUTION
%% ================================================================
\section{The Triangular Distribution}\label{sec:triangular}

\subsection{Statement}\label{sec:tri-statement}

\begin{theorem}[Triangular limit law for per-step displacement]
\label{thm:triangular}
Let $\mu_p$ denote the empirical measure of the displacements
$\{\delta(a/b) : a/b \in \F_{p-1}^*, \, b \ge 2\}$.  Then as
$p \to \infty$ over primes,
\[
  \mu_p \;\Longrightarrow\; \mathrm{Triangular}[-1,1]
\]
in the sense of weak convergence, where $\mathrm{Triangular}[-1,1]$
has density $f(x) = 1 - |x|$ on $[-1,1]$.
\end{theorem}

\begin{remark}[Conditionality]\label{rem:tri-conditional}
Step~1 of the proof below uses an equidistribution estimate for the
pairs $(a/b, \, (pa \bmod b)/b)$ that is uniform in~$p$.  While the
pointwise equidistribution for fixed~$b$ follows from Weyl's criterion
and Ramanujan sum bounds, the \emph{uniformity} over all denominators
$b \le p - 1$ rests on a routine large-sieve type estimate.  We state
the theorem as unconditional because such estimates are standard
(see, e.g., Montgomery--Vaughan~\cite{MonVaughan2007},
Chapter~7), but we note that the explicit error term
$O(b^{-1/3})$ in Step~2 depends on this uniformity.  Supplying a
fully self-contained proof of the uniform bound is a straightforward
but non-trivial exercise that we leave to the interested reader.
\end{remark}

\begin{corollary}[Universal moment formula]\label{cor:moments}
For each fixed $k \ge 1$, the even moment sums satisfy
\begin{equation}\label{eq:moment-formula}
  S_{2k}(p) := \sum_{a/b \in \F_{p-1}^*} \delta(a/b)^{2k}
  = \frac{3}{\pi^2(2k+1)(k+1)}\, p^2 + O_k(p^{5/3}).
\end{equation}
All odd moments vanish: $S_{2k+1}(p) = 0$ for all $k \ge 0$.
\end{corollary}

Special cases:
\begin{center}
\begin{tabular}{lll}
\toprule
$k$ & Predicted $S_{2k}/p^2$ & Empirical ($p = 997$) \\
\midrule
$1$ & $1/(2\pi^2) \approx 0.05066$ & $0.05008$ \\
$2$ & $1/(5\pi^2) \approx 0.02026$ & $0.01984$ \\
$3$ & $3/(28\pi^2) \approx 0.01086$ & $0.01053$ \\
$4$ & $1/(15\pi^2) \approx 0.00675$ & $0.00669$ \\
\bottomrule
\end{tabular}
\end{center}

\subsection{Proof}\label{sec:tri-proof}

The proof proceeds in three steps.

\smallskip
\noindent\textbf{Step 1: Equidistribution for fixed denominator.}
Fix $b \ge 2$.  For $\gcd(a,b) = 1$, consider the pair
$(a/b, \, (pa \bmod b)/b)$ in $[0,1]^2$.  The exponential sum
test gives
\[
  \sum_{\gcd(a,b)=1} e\bigl((h + kp)a/b\bigr) = c_b(h + kp),
\]
where $c_b$ is the Ramanujan sum.  For $(h,k) \ne (0,0)$ and
$b \to \infty$, the bound $|c_b(h + kp)| \le \gcd(h + kp, b) = o(\varphi(b))$
for all but finitely many $b$ (namely, the divisors of $h + kp$).
By the Weyl criterion, the pairs become equidistributed in
$[0,1]^2$ as $b \to \infty$.

Since $\delta(a/b) = a/b - (pa \bmod b)/b$, equidistribution of
the pair $(U, V)$ in $[0,1]^2$ implies $U - V$ converges to
the distribution of $U_1 - U_2$ for independent uniforms $U_i$,
which is $\mathrm{Triangular}[-1,1]$.

\smallskip
\noindent\textbf{Step 2: Averaging over denominators.}
The $2k$-th empirical moment is
\[
  m_{2k}(p) = \frac{S_{2k}(p)}{|\F_{p-1}|}
  = \frac{1}{|\F_{p-1}|}\sum_{b=2}^{p-1} \varphi(b)\, m_{2k}(b,p),
\]
where $m_{2k}(b,p) = \varphi(b)^{-1}\sum_{\gcd(a,b)=1}
\delta(a/b)^{2k}$ is the per-denominator moment.

By Step~1 and the Koksma--Hlawka inequality, for $b > B_0$:
\[
  m_{2k}(b,p) = \frac{2}{(2k+1)(2k+2)} + O(b^{-1/3}).
\]
Splitting at a cutoff $B_0 = p^{5/6}$: the small-$b$ contribution
is $O(B_0^2) = O(p^{5/3})$, and the large-$b$ approximation error
is $O(\sum_{b > B_0} \varphi(b)/b^{1/3}) = O(p^{5/3})$.  Thus
\[
  S_{2k}(p) = \frac{2}{(2k+1)(2k+2)} \cdot \frac{3p^2}{\pi^2}
    + O(p^{5/3})
  = \frac{3}{\pi^2(2k+1)(k+1)}\,p^2 + O(p^{5/3}).
\]

\smallskip
\noindent\textbf{Step 3: Moment determinacy.}
Since $\mathrm{Triangular}[-1,1]$ has compact support, it is
uniquely determined by its moments.  Convergence of all moments
implies weak convergence.

\smallskip
\noindent\textbf{Odd moments vanish:}
The symmetry $\delta((b-a)/b) = -\delta(a/b)$ (from
$p(b-a) \bmod b = b - (pa \bmod b)$) pairs each displacement
with its negative, forcing all odd moments to zero.
\qed

\begin{remark}
The triangular distribution is decisively \emph{not} Gaussian:
the standardized fourth moment (kurtosis) is $12/5 = 2.4$
versus $3$ for a Gaussian.  Empirical values match the triangular
prediction to 3--4 significant figures.
\end{remark}



%% ================================================================
%% SECTION 8: COMPOSITES AND THE HEALING PHENOMENON
%% ================================================================
\section{Composites and the Healing Phenomenon}\label{sec:composites}

\subsection{The dichotomy: primes damage, composites heal}
\label{sec:dichotomy}

The per-step framework reveals a sharp asymmetry:
\begin{itemize}
  \item Primes with $M(p) \le -3$ universally increase discrepancy
    (Sign Theorem).
  \item Composites overwhelmingly decrease it.
\end{itemize}

The balance between damage and healing, governed by the Mertens
function, drives the convergence $W(N) \to 0$.

\subsection{The zero-map mechanism}\label{sec:zero-map}

For composite $N$ and any denominator $b$ dividing $N$, the
multiplication map $\sigma_N \colon a \mapsto Na \bmod b$ sends
every residue to~$0$:
\[
  \delta(a/b) = \frac{a}{b} \quad\text{for all } b \mid N.
\]
This is the \emph{zero-map mechanism}: divisor-denominators
contribute deterministic, structured shifts that are fully
controlled.  Primes have essentially no such protection
(only $b = 1$ divides a prime).

\begin{proposition}\label{prop:composite-advantage}
For a composite $N$ with $d(N)$ divisors, the number of old fractions
with structured (zero-map) shifts is
$\sum_{j=0}^{k-1} \varphi(d_j)$,
where $d_0, d_1, \ldots, d_{k-1}$ are the proper divisors of~$N$.
For even~$N$, this includes all fractions with power-of-2
denominators; for $N$ with $k$ distinct prime factors, this
includes all fractions with smooth denominators.
\end{proposition}

\subsection{The density-zero theorem}\label{sec:density-zero}

\begin{theorem}[Composite healing]\label{thm:composite-heal}
Assume the following asymptotic independence hypothesis:
\begin{quote}
\textbf{(AI)} \emph{For squarefree composites $N = pq$ (semiprimes), the
value $\mu(N) \in \{+1, -1\}$ is asymptotically independent of
$\operatorname{sgn}(M(N-1))$.}
\end{quote}
Then the set $\mathcal{N}_{\mathrm{bad}} = \{N \text{ composite}, \, N \ge 10 :
W(N) \ge W(N-1)\}$ has natural density zero among the composites:
\[
  \frac{\#\{N \le X \text{ composite} : N \in \mathcal{N}_{\mathrm{bad}}\}}
       {\#\{N \le X \text{ composite}\}}
  \;\to\; 0 \quad\text{as } X \to \infty.
\]
\end{theorem}

\begin{remark}
Hypothesis~\textbf{(AI)} asserts that knowing whether a semiprime $pq$
has an even or odd number of prime factors (which is deterministic:
$\mu(pq) = 1$) gives no information about the sign of $M(pq - 1)$.
This is a standard heuristic assumption in analytic number theory,
closely related to Chowla-type conjectures on correlations of
multiplicative functions.  It is used only in Case~(III) of the
proof (the semiprime case).  Cases~(I) and~(II) are unconditional.
\end{remark}

\begin{proof}[Proof sketch]
The argument partitions composites into three classes:

\smallskip
\noindent\textbf{(I) Squarefree composites with $|M(N-1)| \ge (\log N)^2$.}
For such $N$, the Mertens function provides a strong signal: the
four-term decomposition is dominated by the $M$-dependent terms,
and healing occurs when $\mu(N) \cdot M(N-1) < 0$.  Since $\mu(N)$
takes values $\pm 1$ with equal frequency among squarefree integers,
and $M(N-1)$ has definite sign when $|M(N-1)|$ is large, these
composites heal with probability approaching~$1$.  The set where
$|M(N-1)| < (\log N)^2$ has density $O(1/\log N)$ by the oscillation
theory of $M$.

\smallskip
\noindent\textbf{(II) Non-squarefree composites with $\ge 3$ distinct
prime factors.}
By the dilution-dominance argument: $\varphi(N)/N \le \prod_{p \mid N}
(1 - 1/p) < 1$, so the cost of insertion is reduced.  The ratio
$A/(C + D)$ grows with $n/\varphi(N)$, which diverges.  All such
composites heal for $N \ge N_0$.

\smallskip
\noindent\textbf{(III) The remaining composites: semiprimes $pq$ and
prime squares $p^2$.}
Prime squares contribute $O(\sqrt{X})$ integers up to $X$, which is
density zero.  Non-healing semiprimes correspond to those where
$\mu(pq) \cdot M(pq - 1) > 0$, which by the independence of $\mu(pq)$
and $\operatorname{sgn}(M(pq-1))$ constitute at most half the
semiprimes.  Since semiprimes have density $O(1/\log X)$ among
composites, the non-healing density from this class is
$O(1/\log X) \to 0$.

Combining the three classes: the total non-healing density
tends to zero.
\end{proof}

\subsection{Special case: powers of 2}\label{sec:powers-of-2}

\begin{proposition}\label{prop:powers-of-2}
For all $k \ge 3$, $W(2^k) < W(2^{k-1})$.  That is, every sufficiently
large power of~$2$ heals the Farey sequence.
\end{proposition}

\begin{proof}[Proof sketch]
For $N = 2^k$: $\varphi(2^k) = 2^{k-1}$ new fractions are inserted,
and every denominator $b = 2^j$ ($j = 0, \ldots, k-1$) has the
zero-map property.  The dilution ratio $A/(C + D)$ grows as
$n/2^{k-1} \sim 3 \cdot 2^{k-1}/\pi^2$, which diverges.
For $k \ge 3$, the dilution dominates the costs.
(The case $k = 2$, i.e., $N = 4$, is a boundary case where
$\Delta W$ is numerically indistinguishable from zero.)
\end{proof}

\subsection{Non-healing characterization}\label{sec:non-healing-class}

\begin{observation}\label{obs:non-healing}
Computation through $N = 1{,}500$ reveals that non-healing composites
fall into exactly three categories:
\begin{enumerate}[(I)]
  \item \textbf{Prime squares $p^2$} ($p \ge 11$): these consistently
    fail to heal because $\varphi(p^2)/|\F_{p^2-1}|$ falls below
    a critical density threshold.  Count up to $X$: $O(\sqrt{X}/\log X)$.
  \item \textbf{Semiprimes $pq$} with unfavorable Mertens sign:
    $\mu(pq) \cdot M(pq-1) > 0$.  Count up to $X$: $O(X/\log X)$,
    which is $o(X)$ among composites.
  \item \textbf{Sporadic composites} near the Mertens boundary
    $|M(N-1)| \le 2$: density zero by Mertens oscillation.
\end{enumerate}
\end{observation}



%% ================================================================
%% SECTION 9: THE ZERO SPECTRUM
%% ================================================================
\section{The Zero Spectrum}\label{sec:zeros}

\subsection{Explicit formula for $\Delta W(p)$}\label{sec:explicit}

The Mertens function admits the (formal) explicit formula
\[
  M(x) = -\sum_{\rho} \frac{x^{\rho}}{\rho\,\zeta'(\rho)}
    + \text{lower order},
\]
where $\rho$ ranges over the nontrivial zeros of $\zeta(s)$.
Since the bridge identity (Theorem~\ref{thm:bridge}) connects
$\sum e^{2\pi ipf}$ to $M(p)$, and the $L^2$ discrepancy $W(N)$
admits a Parseval-type representation in terms of
$\sigma_m(N) = \sum_{d \mid m} d\,M(\lfloor N/d \rfloor)$,
the per-step change $\Delta W(p)$ inherits a bilinear-form
expansion over pairs of zeta zeros:
\begin{equation}\label{eq:explicit-bilinear}
  \Delta W(p) \sim \frac{\pi^2}{12p^2}\Bigl[
    \sum_{\gamma} |A(\rho)|^2 K(\rho,\rho)\,p
    + 2\operatorname{Re}\!\sum_{\gamma < \gamma'} S(\gamma,\gamma';p)\,p
    - 2\operatorname{Re}\!\sum_{\gamma} A(\rho)\,p^{\rho}\,T_1(\rho,p)
  \Bigr],
\end{equation}
where $A(\rho) = -1/(\rho\,\zeta'(\rho))$, and
$K(\rho,\rho')$, $S$, $T_1$ are explicit arithmetic functions.

The three terms have transparent interpretations:
\begin{enumerate}[(i)]
  \item \textbf{Diagonal} ($\gamma = \gamma'$): always positive, the
    ``healing pressure'' from statistical averaging.
  \item \textbf{Off-diagonal oscillation}: oscillatory at beat
    frequencies $\gamma - \gamma'$, creating the quasi-random
    pattern in $\Delta W(p)$.
  \item \textbf{$M(p)$ coupling}: proportional to $-M(p)$, the
    mechanism by which the Mertens function influences healing.
\end{enumerate}

\begin{remark}
The derivation~\eqref{eq:explicit-bilinear} is formal: making the
error terms rigorous requires controlling the truncation of the
zero sum and the regularization of the diagonal.  The qualitative
structure is robust and matches all numerical evidence.
\end{remark}

\subsection{Detection of zeta zeros in the Farey spectrum}
\label{sec:detection}

The explicit formula predicts that the Fourier transform of
$\Delta W(p) \cdot p^2$ (sampled at $t = \log p$ over primes) should
exhibit peaks at the imaginary parts $\gamma_k$ of zeta zeros.

\begin{observation}[Zero detection]\label{obs:zero-detection}
Lomb--Scargle periodogram analysis of $\Delta W(p) \cdot p^2$ for
9,588 primes up to $10^5$ confirms this prediction with overwhelming
statistical significance:
\begin{center}
\begin{tabular}{lrr}
\toprule
Frequency & $Z$-score & $p$-value \\
\midrule
$\gamma_1 = 14.135$ & $1{,}625$ & $< 0.001$ \\
$\gamma_2 = 21.022$ & $470$ & $< 0.001$ \\
$\gamma_3 = 25.011$ & $886$ & $< 0.001$ \\
$\gamma_4 = 30.425$ & $446$ & $< 0.001$ \\
$\gamma_5 = 32.935$ & $91$ & $< 0.001$ \\
$\gamma_2 - \gamma_1 = 6.887$ & $85$ & $< 0.001$ \\
$\gamma_3 - \gamma_1 = 10.876$ & $251$ & $< 0.001$ \\
$2\gamma_1 = 28.269$ & $167$ & $< 0.001$ \\
\bottomrule
\end{tabular}
\end{center}
$Z$-scores are computed against 100 random reshufflings of $\Delta W$
values with $\log p$ positions fixed.  All tested frequencies are
detected; the two strongest peaks in the full periodogram correspond
to the first two zeta zeros.
\end{observation}

\subsection{Interpretation}\label{sec:zeros-interpret}

This detection is \emph{expected} from the explicit formula --- it
confirms the theory rather than discovering new information about
zeros.  The hierarchy of signal strengths is:
\begin{enumerate}
  \item Single zeros $\gamma_k$: very strong;
  \item Differences $\gamma_k - \gamma_l$: moderate, requiring more
    data for small differences;
  \item Harmonics $2\gamma_1$: present and significant.
\end{enumerate}

The signal \emph{strengthens} with more primes (rather than washing
out), confirming it is a genuine spectral feature.  However, this
does \textbf{not} provide new information about the location of
zeros: detecting zeta-zero content in $\Delta W$ is a filtered
version of detecting it in $M(x)$, which is the standard approach
and carries more information.



%% ================================================================
%% SECTION 10: FORMAL VERIFICATION
%% ================================================================
\section{Formal Verification}\label{sec:lean}

\subsection{What is verified}\label{sec:lean-verified}

All core identities in this paper are formally verified in
Lean~4 using the Mathlib library.  The verification encompasses
258 results across fifteen files, including:

\begin{itemize}
  \item The bridge identity (Theorem~\ref{thm:bridge}) and its
    generalization (Theorem~\ref{thm:bridge}, eq.~\eqref{eq:universal}).
  \item The four-term decomposition (Theorem~\ref{thm:four-term}).
  \item The permutation square-sum identity (Theorem~\ref{thm:perm-sq}).
  \item The displacement--shift identity (Proposition~\ref{prop:disp-shift}).
  \item The Dedekind sum identity (Lemma~\ref{lem:dedekind-identity}).
  \item The M\"obius reduction (Lemma~\ref{lem:mobius-reduction}).
  \item The spectral factorization (Theorem~\ref{thm:spectral}).
  \item The Sign Theorem for specific small primes ($p = 13, 19, 31$,
    verified by \texttt{native\_decide}).
\end{itemize}

The bridge-identity proof chain was constructed autonomously by the
Aristotle automated theorem prover (an API-based tool for Lean~4
proof search).

\subsection{What is not verified}\label{sec:lean-gaps}

One \texttt{sorry} remains in the Lean formalization:
\texttt{sign\_theorem\_conj}, the general statement of the Sign
Theorem for all primes with $M(p) \le -3$.  The analytical proof
(Section~\ref{sec:sign}) relies on Dedekind reciprocity and
asymptotic estimates that have not yet been formalized in Lean~4.

The Dedekind reciprocity law (Theorem~\ref{thm:reciprocity}) is
used in the proof but has not been formalized in our codebase.
(A Lean formalization exists in the Mathlib library for the
closely related theory of modular forms, but we have not yet
connected it to our framework.)

\subsection{The Aristotle theorem prover}\label{sec:aristotle}

Aristotle is an automated theorem prover for Lean~4, accessed via
the \texttt{harmonic.fun} API.  It was used to:
\begin{itemize}
  \item Search for proof tactics for lemmas involving Ramanujan sums,
    M\"obius inversion, and Farey sequence combinatorics.
  \item Close auxiliary lemmas in the bridge-identity proof chain.
  \item Verify the four-term decomposition identity by symbolic
    expansion.
\end{itemize}

\subsection{Reproducibility}\label{sec:reproducibility}

All computational results (wobble values, Mertens function
evaluations, spectral analysis) are reproducible from the scripts
in the \texttt{experiments/} directory of the project repository.
Key computations use exact rational arithmetic (Python
\texttt{fractions} module) or high-precision floating point
(256-bit MPFR via \texttt{mpmath}) for the critical counterexample
at $p = 92{,}173$.



%% ================================================================
%% SECTION 11: OPEN QUESTIONS
%% ================================================================
\section{Open Questions}\label{sec:open}

\begin{enumerate}
  \item \textbf{Threshold improvement.}
    Can the Sign Theorem be extended to $M(p) \le -2$?
    The unique counterexample at $p = 92{,}173$ has $M(p) = -2$;
    is this the only one?  More ambitiously: for which function
    $g(p)$ does $M(p) \le -g(p)$ guarantee $\Delta W(p) < 0$?

  \item \textbf{$B \ge 0$ conjecture.}
    Prove or disprove: $B(p) \ge 0$ for all primes $p \ge 11$
    with $M(p) \le -3$.  This is verified computationally for
    all qualifying primes up to $10^5$, but no analytical proof
    is known.

  \item \textbf{$C_W = O(1)$ conjecture.}
    Prove that $N \cdot W(N)$ is bounded by an absolute constant.
    This is weaker than RH (which gives $N W(N) = O(N^{\varepsilon})$)
    but stronger than the unconditional bound $C_W \le \log N$.
    Empirically, $C_W$ appears to converge to approximately $0.67$.

  \item \textbf{Triangular error term.}
    What is the optimal error in Corollary~\ref{cor:moments}?
    The current bound $O(p^{5/3})$ comes from a crude discrepancy
    estimate; sharper bounds (possibly $O(p^{3/2+\varepsilon})$)
    may follow from Kloosterman-sum estimates.

  \item \textbf{Spectral energy lower bound.}
    Prove that $\sum_{\chi \text{ odd}} |\Lambda_p(\chi)|^2
    \gg p^2 \log p$, where $\Lambda_p(\chi)$ is the character
    transform of the Mertens-weighted function
    $\lambda_p(a) = M(\lfloor (p-1)/a \rfloor) + \mathbf{1}_{a=1}$.
    This would connect the spectral positivity of
    Corollary~\ref{cor:spectral-pos} to quantitative lower bounds
    on $\Delta W$.

  \item \textbf{Other discrepancy measures.}
    Does the primes-damage-composites-heal dichotomy persist for
    the $L^{\infty}$ discrepancy, the star discrepancy, or
    entropy-based measures?

  \item \textbf{Higher dimensions.}
    For the tensor-product Farey set
    $\F_N^d = \F_N \times \cdots \times \F_N \subset [0,1]^d$,
    the bridge identity gives
    $\sum_{f \in \F_N^d} e^{2\pi i m \cdot f} = (M(N)+1)^d$.
    Does the per-step framework generalize to dimension $d \ge 2$,
    and does the primes/composites dichotomy persist?
\end{enumerate}



%% ================================================================
%% SECTION 12: BIBLIOGRAPHY
%% ================================================================
\begin{thebibliography}{99}

\bibitem{Franel1924}
J.~Franel,
Les suites de Farey et le probl\`eme des nombres premiers,
\textit{G\"ottinger Nachrichten} (1924), 198--201.

\bibitem{Landau1924}
E.~Landau,
Bemerkungen zu der vorstehenden Abhandlung von Herrn Franel,
\textit{G\"ottinger Nachrichten} (1924), 202--206.

\bibitem{HardyWright2008}
G.\,H.~Hardy and E.\,M.~Wright,
\textit{An Introduction to the Theory of Numbers},
6th ed., Oxford University Press, 2008.

\bibitem{Edwards1974}
H.\,M.~Edwards,
\textit{Riemann's Zeta Function},
Academic Press, 1974; reprinted Dover, 2001.

\bibitem{Ramanujan1918}
S.~Ramanujan,
On certain trigonometrical sums and their applications
in the theory of numbers,
\textit{Trans.\ Cambridge Philos.\ Soc.}
\textbf{22} (1918), no.~13, 259--276.

\bibitem{Kanemitsu1982}
S.~Kanemitsu and R.~Sita~Rama~Chandra~Rao,
On a conjecture of S.~Chowla and of S.~Chowla and H.~Walum, II,
\textit{J.~Number Theory} \textbf{20} (1985), 103--120.

\bibitem{RademacherGrosswald1972}
H.~Rademacher and E.~Grosswald,
\textit{Dedekind Sums},
Carus Mathematical Monographs~16, MAA, 1972.

\bibitem{Walfisz1963}
A.~Walfisz,
\textit{Weylsche Exponentialsummen in der neueren Zahlentheorie},
VEB Deutscher Verlag der Wissenschaften, Berlin, 1963.

\bibitem{ElMarraki1995}
M.~El~Marraki,
Fonction sommatoire de la fonction de M\"obius, III:
Majorations asymptotiques effectives fortes,
\textit{J.~Th\'eor.~Nombres Bordeaux} \textbf{7} (1995), 407--433.

\bibitem{Montgomery1973}
H.\,L.~Montgomery,
The pair correlation of zeros of the zeta function,
\textit{Proc.\ Sympos.\ Pure Math.}
\textbf{24} (1973), 181--193.

\bibitem{Vardi1987}
I.~Vardi,
A relation between Dedekind sums and Kloosterman sums,
\textit{Duke Math.~J.} \textbf{55} (1987), no.~1, 189--197.

\bibitem{ErdosTuran1948}
P.~Erd\H{o}s and P.~Tur\'an,
On a problem in the theory of uniform distribution, I--II,
\textit{Indag.\ Math.} \textbf{10} (1948), 370--378, 406--413.

\bibitem{RubinsteinSarnak1994}
M.~Rubinstein and P.~Sarnak,
Chebyshev's bias,
\textit{Experiment.\ Math.} \textbf{3} (1994), no.~3, 173--197.

\bibitem{KarvonenZhigljavsky2025}
A.~Karvonen and A.~Zhigljavsky,
Maximum mean discrepancy of Farey sequences,
\textit{J.~Number Theory} (2025), to appear.

\bibitem{Goldstein1973}
L.\,J.~Goldstein,
Dedekind sums for a Fuchsian group, I,
\textit{Nagoya Math.~J.} \textbf{50} (1973), 21--47.

\bibitem{MonVaughan2007}
H.\,L.~Montgomery and R.\,C.~Vaughan,
\textit{Multiplicative Number Theory I: Classical Theory},
Cambridge University Press, 2007.

\bibitem{Titchmarsh1986}
E.\,C.~Titchmarsh,
\textit{The Theory of the Riemann Zeta-Function},
2nd ed.\ (revised by D.\,R.~Heath-Brown),
Oxford University Press, 1986.

\bibitem{IwaniecKowalski2004}
H.~Iwaniec and E.~Kowalski,
\textit{Analytic Number Theory},
AMS Colloquium Publications~53, 2004.

\bibitem{Niederreiter1992}
H.~Niederreiter,
\textit{Random Number Generation and Quasi-Monte Carlo Methods},
SIAM, 1992.

\bibitem{Tenenbaum2015}
G.~Tenenbaum,
\textit{Introduction to Analytic and Probabilistic Number Theory},
3rd ed., AMS, 2015.

\bibitem{DressMendesFrance1988}
F.~Dress and M.~Mend\`es~France,
On the distribution of the Farey sequence,
\textit{J.~Number Theory} \textbf{29} (1988), 244--252.

\bibitem{Ingham1942}
A.\,E.~Ingham,
On two conjectures in the theory of numbers,
\textit{Amer.~J.~Math.} \textbf{64} (1942), 313--319.

\bibitem{Ng2004}
N.~Ng,
The distribution of the summatory function of the M\"obius function,
\textit{Proc.\ London Math.\ Soc.} \textbf{89} (2004), 361--389.

\bibitem{Goldston1984}
D.\,A.~Goldston,
The second moment for prime numbers,
\textit{Quart.~J.~Math.\ Oxford} \textbf{35} (1984), 153--163.

\bibitem{Rademacher1938}
H.~Rademacher,
On the partition function $p(n)$,
\textit{Proc.\ London Math.\ Soc.} \textbf{43} (1938), 241--254.

\bibitem{Myerson1989}
G.~Myerson,
Dedekind sums and uniform distribution,
\textit{J.~Number Theory} \textbf{28} (1988), no.~3, 233--239.

\end{thebibliography}


\end{document}
